

\section{مقدمة}

بعد الانتهاء من بناء النموذج الحركي والديناميكي للذراع الروبوتية في الفصول السابقة، يتم في هذا الفصل الانتقال إلى مرحلة تصميم نظام التحكم. تمثل هذه المرحلة حلقة الوصل بين النموذج الرياضي المستنتج وبين السلوك المطلوب من الذراع أثناء التشغيل، إذ لا يقتصر الهدف على توصيف حركة الذراع فقط، بل يمتد إلى توليد إشارات تحكم مناسبة للمشغلات بحيث تستطيع المفاصل تتبع المواضع أو المسارات المرجعية المطلوبة.

يعتمد تصميم نظام التحكم على وجود نموذج رياضي قادر على توصيف العلاقة بين دخل المشغلات والحركة الناتجة عن مفاصل الذراع. لذلك فإن النموذج الديناميكي المستنتج سابقًا لا يستخدم فقط في تحليل العزوم والقوى المؤثرة في الذراع، بل يشكل أيضًا الأساس الذي يتم من خلاله بناء المتحكم واختبار استجابة النظام في بيئة المحاكاة.

قبل الدخول في تصميم المتحكمات، لا بد من إعادة صياغة النموذج الديناميكي بالشكل المناسب لأغراض التحكم. لذلك يبدأ هذا الفصل بتجهيز النموذج الرياضي، وذلك من خلال عزل التسارع من معادلة الحركة، ثم دمج نموذج المشغلات مع النموذج الديناميكي للذراع، وبعد ذلك تحويل النموذج الناتج إلى تمثيل فضاء حالة غير خطي. تسمح هذه الصياغة بتمثيل النظام على الصورة العامة:
\[
\dot{x}=f(x,u)
\]
حيث يمثل \(x\) متجه الحالة، ويمثل \(u\) دخل التحكم المطبق على المشغلات.

بعد تجهيز النموذج، يتم توضيح البنية العامة لنظام التحكم المعتمد في هذا العمل. تقسم هذه البنية إلى مستويين رئيسيين: المستوى الأول هو التحكم منخفض المستوى \LR{Low-Level Control}، وهو المسؤول عن التحكم بمواضع المحركات أو المفاصل وتتبع القيم المرجعية المطلوبة. أما المستوى الثاني فهو التحكم عالي المستوى \LR{High-Level Control}، وهو المسؤول عن التعامل مع حركة الذراع ككل، مثل توليد الأوامر المرجعية أو تحديد الحركة المطلوبة للنهاية المنفذة ثم تمريرها إلى مستوى التحكم بالمفاصل.

وبناءً على ذلك، يعرض هذا الفصل أولًا خطوات تجهيز النموذج الديناميكي للتحكم، ثم يوضح بنية التحكم المعتمدة، وبعدها يتم عرض تصميم المتحكم منخفض المستوى الخاص بمواضع المفاصل، ثم المتحكم عالي المستوى المسؤول عن حركة الذراع ككل.